\documentclass[10pt,a4paper]{amsart}
\usepackage[margin=19mm]{geometry}
\usepackage{amsmath,amssymb,mathtools}
\usepackage[scr=rsfs,cal=euler]{mathalfa}
\usepackage{xfrac}
\usepackage[unicode,hidelinks,bookmarks=false]{hyperref}
\newcommand{\ie}{i.e.\ }
\newcommand{\cf}{cf.\ }
\newcommand{\resp}{respectively\ }
\newcommand{\Subsection}{\subsection}
\providecommand{\nobreakdash}{\nobreak\hbox{-}}
\setlength{\emergencystretch}{2em}
\multlinegap0pt
\usepackage{amssymb}
\usepackage{enumerate}
\newtheorem{lem}{Lemma}
\newcommand{\D}{\mathbb{D}}
\newcommand{\abs}[1]{\left\vert#1\right\vert}
\newcommand{\hot}{{\mathrm{h.o.t.}}}
\newcommand{\inv}{^{-1}}
\newcommand{\rest}[1]{ \arrowvert_{#1}}
\title{Degenerate homoclinic bifurcations in complex dimension 2}
\author{Romain Dujardin}
\date{Source: arXiv:2306.08160v3 (CC BY 4.0)}
\begin{document}
\maketitle

\noindent\emph{Source and licence.} Degenerate homoclinic bifurcations in complex dimension 2. Version arXiv:2306.08160v3 (CC BY 4.0), distributed by arXiv under the Creative Commons Attribution 4.0 International licence, http://creativecommons.org/licenses/by/4.0/, as recorded in the arXiv licence metadata for this exact version and shown on the abstract page https://arxiv.org/abs/2306.08160v3. Full licence notice: \texttt{licenses/arxiv-2306.08160-CC-BY-4.0.txt}.\par\smallskip

\noindent\textit{Editorial scope.} Counted unit: the article's lem lem:elementary\_analysis (source0.tex lines 1856--1864). The article's complete original proof of it is delivered with it (source0.tex lines 1866--1903, one \texttt{proof} environment of 38 lines). No mathematics is written, repaired or re-derived: the statement and the proof are the article's own lines, in the article's own notation, and no retained line is substituted or re-flowed.  No further source-local context is needed: every cross-reference of every retained line resolves inside the two delivered spans.  Nothing was omitted from any retained span, so the package carries no omission record.  No retained line contains a citation, so the wrapper carries no bibliography and no bibliography entry.  No retained line contains a figure environment, an image-inclusion command or drawing code, so no image is delivered and the package declares no asset dependency.  Editorial formatting only, in addition: an AMS \texttt{amsart} preamble that restates the article's own declarations and macros used by the retained lines, namely the environments \texttt{lem} on separate counters, together with the standard packages those lines need.  The retained environments print in this document as Lemma 1 in order of appearance, whereas the article prints the same items with the numbering of its own full document.  The complete native source of the article as distributed in the arXiv e-print for this exact version is bundled unchanged as \texttt{sources/142892/source0.tex}.
\begin{lem}\label{lem:elementary_analysis}
If  $ \delta < \abs{u_0}\inv$   is fixed, then  for 
  sufficiently large $n$,  there are $m$ disjoint topological disks $\Lambda_{n, i}\subset \Lambda$, 
  $1\leq i\leq m$, with $\Lambda_{n, i} \subset D(0, C \abs{u_0}^{-n/m})$, 
  in which 
$\lambda \mapsto x(\lambda ) - \alpha(\lambda) u_\lambda^{-n} $ realizes a 
biholomorphism
$\Lambda_{n, i}\to D(0, \delta^n)$. 
\end{lem}

\begin{proof}[Proof of Lemma \ref{lem:elementary_analysis}] We  rely on the following   fact 
 from elementary  complex analysis, which we leave as an exercise to the reader: if $f$ is a holomorphic function on $\D$ such that $f(0)  = 0$, $f'(0) = 1$ and $\abs{f'}\leq M$, then there exists $\rho= \rho(M)$ 
and a domain $\Omega$ with $D(0, \rho/2)\subset \Omega\subset D(0, 3\rho/2)$ such that $f\rest{\omega}:\Omega\to D(0, \rho)$ is a biholomorphism. By rescaling it holds with $f$ holomorphic in $D(0, R)$, 
$\abs{f'(0)} = d$, $\abs{f'}\leq dM$ and the image radius is $\rho Rd$. 


Recall that by assumption $x(\lambda) = \lambda^m+\hot$
 Let us first show that there are $m$ solutions to the equation $x(\lambda)   = \alpha(\lambda) u_\lambda^{-n}$ close to the origin.
 Write  $x(\lambda) = g(\lambda)^m$, where $g$ is holomorphic in some disk 
$D(0, R)$ and $g'(0) = 1$, so that  by the above result 
$g$ is a univalent  map  $\Omega \to D(0,  \rho R)$, for some domain $\Omega$.

Recall that if 
 $\lambda$ is exponentially small (i.e. $\abs{\lambda}\leq (1-\eta)^n$ for some $\eta>0$), then
$u_\lambda^{-n}\sim  u_0^{-n}$: indeed $u_\lambda = u(1+O(\lambda))$, so 
\begin{align}\label{eq:ulambda}
u_\lambda^{-n} &= u_0^{-n}(1+ O((1-\eta)^n))^n = u_0^{-n} \exp(-n\ln(1+O((1-\eta)^n) ))   \\ & \notag
= u_0^{-n} \exp(n O((1-\eta)^n) )  
  \sim u_0^{-n}. 
\end{align}
 It follows that for every choice $\zeta_i$, $1\leq i \leq m$ 
  of  $m^{\rm th}$ root of $\alpha(\lambda) u_\lambda^{-n}$, we have a solution $\lambda_{n,i}$ 
of $g(\lambda)   = \zeta_i$ in $\Omega$ with $\lambda_{n,i}\asymp \abs{u_0}^{-n/m}$. 
%(For notational convenience we drop the $i$ and write $\lambda_n = \lambda_{n,i}$.)
At $\lambda_{n,i}$ we have $x'(\lambda_{n,i}) 
\sim m \abs{\lambda_{n,i}}^{m-1} \asymp
\abs{u_0}^{-n({m-1})/{m}}$, while reasoning as in \eqref{eq:ulambda} we get that 
$\frac{d}{d\lambda}\big\vert_{\lambda = \lambda_{n, i}}( \alpha(\lambda) u_\lambda^{-n}) 
= O(\abs{u_0}^{-n})$. Let $\beta$ be such that 
$ \delta\abs{u_0}^{(m-1)/m}< \beta < u^{-1/m}$. Then in $D(\lambda_{n,i}, \beta^n)$ we get that 
$\abs{\frac{d}{d\lambda}(x(\lambda) - \alpha(\lambda) u_\lambda^{-n}} \leq C \abs{u_0}^{-n({m-1})/{m}}$, 
so by the preliminary fact, 
$\lambda\mapsto x(\lambda) - \alpha(\lambda) u_\lambda^{-n}$ realizes  a biholomorphism 
 from an approximately round domain $\Lambda_n$
 of size $\asymp \beta^n$  about $\lambda_{n,i}$
 to a disk centered at the origin and 
 of radius $\asymp  \beta^n\abs{u_0}^{- n(m-1)/m}\gg \delta^n$, and we are done. 
\end{proof}

\end{document}
